% GDZ PPN236018647,native498;HKUST499;original diary4r.
% Entries corresponding to printed59-62; final entry63 excluded.
\clearpage\SourceLeaf{498}\NoSourceHead
\begingroup\fontsize{10}{15}\selectfont\itshape
\newcommand{\DiaryStrike}[1]{\tikz[baseline=(t.base)]{\node[inner sep=0pt,outer sep=0pt] (t) {\strut #1};\draw[line width=.55pt] ([yshift=.15ex]t.west)--([yshift=.15ex]t.east);}}
\noindent Formularum integralium formae
\[
\int e^{-t^a}dt\qquad\text{et}\qquad\int\frac{du}{{}^{\beta}\!\surd1+u^\gamma)}
\]
\noindent inter se comparationem institui.\hfill Mrt 2.

\medskip
\noindent Cur \DiaryStrike{plures} ad aequationem perveniatur\\
\hspace*{2em}gradus \(nn^{\underline{\underline{ti}}}\) dividendo curvam lemniscatam\\
\hspace*{2em}in \(n\) partes\quad.\kern.4em.\kern.4em\raisebox{.25ex}{\rule{.45em}{.3pt}}\kern.4em\raisebox{.25ex}{\rule{.7em}{.3pt}}\hfill Mart. 19

\medskip
\noindent A potestatibus Integr. \(\displaystyle\int\frac{dx}{\surd(1-x^4)}\)\raisebox{1.2ex}{\((0-1)\)} pendet
\[
\Sigma\left(\frac{mm+6mn+nn}{(mm+nn)^4}\right)^k
\]

\medskip
\noindent\tikz[baseline=-.4ex,x=1em,y=1em]{
\draw[line width=.25pt] (0,-.65)--(.27,1.05);
\draw[line width=.25pt] (.18,-.65)--(.45,1.05);
\draw[line width=.25pt] (.37,-.65)--(.64,.95);
\draw[line width=.7pt] (-.17,.6)--(.8,.53);
\draw[line width=.7pt] (-.17,.32)--(.8,.25);
\draw[line width=.7pt] (-.17,-.12)--(.8,-.13);
}\quad Lemniscata geometrice in quinque\\
\hspace*{2em}partes dividitur\hfill Mrt. 21.
\endgroup
